\documentclass[11pt,a4paper]{article}
\usepackage[T1]{fontenc}
\usepackage[utf8]{inputenc}
\usepackage[ngerman]{babel}
\usepackage[margin=2.5cm]{geometry}
\usepackage[hidelinks]{hyperref}
\setlength{\parindent}{0pt}
\begin{document}
\begin{center}
  {\Large\textbf{Webtechnologien -- Übungsblatt 1}}
\end{center}
\section*{Aufgabe 5: HTML-Tabelle}
\textbf{Aufgabe:} Ausgangsdatei untersuchen und die Tabelle nach Vorlage
ergänzen. Uni-Links und Quellenlink hinzufügen. CSS unverändert lassen.

\medskip
\textbf{Skizze der Ausgangsseite:} Überschrift
\glqq Top 5 universities in the world\grqq{} und eine Tabellenzeile:
\begin{center}
\begin{tabular}{|l|l|l|}
\hline
\textbf{University of Oxford} & United Kingdom & 98.5 \\ \hline
\end{tabular}
\end{center}
\textbf{CSS-Wirkung:} Rahmen und Innenabstände, farbige Spaltenköpfe,
abwechselnd helle Zeilen und Beschriftung unter der Tabelle.
Zeilenköpfe sind linksbündig; Zahlen und Länder stehen zentriert.

\medskip
\textbf{Ergebnis:} \texttt{ex5.html} enthält fünf Universitäten mit
Land, Ranking und Webseiten-Link. Spaltenköpfe verwenden
\texttt{scope="col"}, Zeilenköpfe \texttt{scope="row"}.
\texttt{thead} und \texttt{tbody} gliedern die Tabelle.
Die Beschriftung \glqq University rankings 2025\grqq{} verlinkt die
vorgegebene Quelle. Die CSS-Regeln sind unverändert.

\medskip
\textbf{Test:} Über Apache unter
\url{http://localhost/webtech-aufgabe5/ex5.html} erreichbar,
Status \texttt{200 OK}.
\medskip
{\small \textbf{Hilfsmittel:} Codex für Erstellung, Prüfung und Dokumentation.}
\end{document}
